% Problema: límits de quocients (valencià).

\begin{exercise}[Límits de quocients]
Calcula els límits següents:
\[
  \text{(i)}\ \lim_{x \to 2} \frac{x^2 - 4}{x - 2} ,
  \qquad
  \text{(ii)}\ \lim_{x \to 1} \frac{x^3 - 1}{x^2 - 1} ,
  \qquad
  \text{(iii)}\ \lim_{x \to 0} \frac{\sqrt{1 + x} - 1}{x} .
\]

\dmarks{3}

\begin{hint}
En els tres, el numerador i el denominador s'anul·len en el punt. Factoritza
per a simplificar el factor comú; en l'últim, multiplica i divideix per
$\sqrt{1 + x} + 1$.
\end{hint}

\begin{answer}
(i) $4$; (ii) $3/2$; (iii) $1/2$.
\end{answer}

\begin{solution}
Com que el límit no depén del valor en el punt, es pot simplificar per a $x$
diferent d'ell i avaluar després.
\begin{parts}
  \item Per a $x \neq 2$, $\dfrac{x^2 - 4}{x - 2} = x + 2$, així que el
        límit és $4$.
  \item Per a $x \neq \pm 1$,
        \[
          \frac{x^3 - 1}{x^2 - 1}
          = \frac{(x - 1)(x^2 + x + 1)}{(x - 1)(x + 1)}
          = \frac{x^2 + x + 1}{x + 1} ,
        \]
        que tendeix a $3/2$.
  \item Per a $x \neq 0$,
        \[
          \frac{\sqrt{1 + x} - 1}{x}
          = \frac{(1 + x) - 1}{x\,\bigl(\sqrt{1 + x} + 1\bigr)}
          = \frac{1}{\sqrt{1 + x} + 1} ,
        \]
        que tendeix a $1/2$.
\end{parts}
\end{solution}

\begin{marking}
Un punt per apartat. Escriure «$0/0$» com a resultat no puntua; simplificar
bé i equivocar-se en avaluar, mig punt.
\end{marking}
\end{exercise}
